\chapter{Implementação Computacional da GAT em Python}
\label{ch:implementacao-computacional}

Este capítulo fornece a implementação computacional completa dos conceitos centrais da Teoria Geométrica de Arbitragem (GAT). Todos os códigos são funcionais em Python 3.12+ com dependências exclusivas em \texttt{numpy} e \texttt{scipy}. O objetivo é que o estudante possa executar, modificar e experimentar cada algoritmo, consolidando a intuição geométrica desenvolvida nos capítulos anteriores.

\section{Simulação de Movimento Browniano Geométrico}
\label{sec:simulacao-gbm}

O ponto de partida é simular a dinâmica estocástica de ativos financeiros sob a hipótese de mercado eficiente (ausência de arbitragem). O Movimento Browniano Geométrico (GBM) é governado pela equação diferencial estocástica:
\begin{equation}
    dS_t = \mu S_t\,dt + \sigma S_t\,dW_t,
\end{equation}
onde \(\mu\) é o drift (retorno esperado), \(\sigma\) é a volatilidade e \(dW_t \sim \mathcal{N}(0, dt)\) é o incremento de Wiener.

A discretização de Euler-Maruyama para um passo \(\Delta t\) é:
\begin{equation}
    S_{t+\Delta t} = S_t \exp\!\left[\left(\mu - \frac{\sigma^2}{2}\right)\Delta t + \sigma\sqrt{\Delta t}\,\varepsilon_t\right], \quad \varepsilon_t \sim \mathcal{N}(0,1).
\end{equation}
A forma exponencial garante positividade exata de \(S_t\), evitando valores negativos que surgiriam na discretização aditiva ingênua.

\begin{lstlisting}[language=Python, caption={Simulação de $M=10$ trajetórias de GBM com $S_0=100$, $\mu=0.08$, $\sigma=0.25$, $T=1$ ano, $N=252$ passos.}]
import numpy as np
import matplotlib.pyplot as plt

def simular_gbm(S0, mu, sigma, T, N, M, seed=42):
    """
    Simula M trajetorias de Movimento Browniano Geometrico.

    Parametros
    ----------
    S0 : float  -- preco inicial
    mu : float  -- drift anualizado
    sigma : float -- volatilidade anualizada
    T  : float  -- horizonte temporal (anos)
    N  : int    -- numero de passos temporais
    M  : int    -- numero de trajetorias
    seed : int  -- semente aleatoria

    Retorna
    -------
    t : ndarray (N+1,)  -- grade temporal
    S : ndarray (N+1, M) -- precos simulados
    """
    rng = np.random.default_rng(seed)
    dt = T / N
    t = np.linspace(0, T, N + 1)
    S = np.zeros((N + 1, M))
    S[0, :] = S0
    for i in range(N):
        eps = rng.normal(0, 1, M)
        S[i + 1, :] = S[i, :] * np.exp(
            (mu - 0.5 * sigma**2) * dt + sigma * np.sqrt(dt) * eps
        )
    return t, S

# Parametros
S0, mu, sigma = 100.0, 0.08, 0.25
T, N, M = 1.0, 252, 10

t, S = simular_gbm(S0, mu, sigma, T, N, M)

fig, ax = plt.subplots(figsize=(10, 5))
for j in range(M):
    ax.plot(t, S[:, j], alpha=0.7, lw=1.2)
ax.axhline(S0, color='gray', ls='--', alpha=0.5)
ax.set_xlabel('Tempo (anos)')
ax.set_ylabel('Preco')
ax.set_title(f'GBM: $S_0={S0}$, $\\mu={mu}$, $\\sigma={sigma}$, $T={T}$')
ax.grid(True, alpha=0.3)
plt.tight_layout()
plt.savefig('gbm_trajetorias.pdf')
plt.show()
\end{lstlisting}

O estudante deve observar que, para \(\mu = r\) (taxa livre de risco), o processo descontado \(\tilde{S}_t = e^{-rt}S_t\) é um martingale --- propriedade que será violada quando existir arbitragem.

\section{Cálculo Numérico da Curvatura}
\label{sec:calculo-curvatura}

\subsection{Fórmula do Tensor de Curvatura}

No contexto da GAT, a conexão de gauge \(A_\mu\) encapsula os excessos de retorno em relação à taxa livre de risco. Para um mercado com \(N\) ativos, a curvatura (forma de curvatura 2) é dada por:
\begin{equation}
    R_{\mu\nu} = \partial_\mu A_\nu - \partial_\nu A_\mu + [A_\mu, A_\nu],
\end{equation}
onde o comutador \([A_\mu, A_\nu] = A_\mu A_\nu - A_\nu A_\mu\) é não nulo apenas para teorias de gauge não-abelianas. No caso abeliano (um único fator de risco), a curvatura reduz-se a \(R_{\mu\nu} = \partial_\mu A_\nu - \partial_\nu A_\mu\).

Para o mercado financeiro, identificamos os índices \(\mu = j\) (ativo) e \(\nu = t\) (tempo), de modo que:
\begin{equation}
    R_{jt} = \frac{\partial A_j}{\partial t} - \frac{\partial A_t}{\partial S_j} + [A_j, A_t].
\end{equation}
Em coordenadas de preço, com \(A_j\) proporcional ao drift ajustado \(\mu_j - r\), a curvatura mede o ``rotacional'' do campo de gauge no espaço \((S, t)\).

\subsection{Implementação por Diferenças Finitas}

Para calcular \(R\) numericamente, aproximamos as derivadas parciais por diferenças finitas centradas de segunda ordem:
\begin{align}
    \frac{\partial A}{\partial S}\bigg|_{(j,k)} &\approx \frac{A_{j+1,k} - A_{j-1,k}}{2\,\Delta S}, \\
    \frac{\partial A}{\partial t}\bigg|_{(j,k)}  &\approx \frac{A_{j,k+1} - A_{j,k-1}}{2\,\Delta t}.
\end{align}

\begin{lstlisting}[language=Python, caption={Cálculo numérico da curvatura $R$ via diferenças finitas. Verificação: $R \approx 0$ quando $\mu_j = r$ (ausência de arbitragem).}]
import numpy as np
from scipy.ndimage import sobel

def calcular_curvatura(precos, mu, r, sigma, dt, dS):
    """
    Calcula o tensor de curvatura R para um mercado com N ativos.

    Parametros
    ----------
    precos : ndarray (N_t, N_assets) -- precos simulados
    mu     : ndarray (N_assets,)     -- drifts
    r      : float                    -- taxa livre de risco
    sigma  : ndarray (N_assets,)     -- volatilidades
    dt     : float                    -- passo temporal
    dS     : float                    -- passo espacial (preco)

    Retorna
    -------
    R : ndarray (N_t-2, N_assets-2) -- curvatura R_{j,t}
    """
    N_t, N_assets = precos.shape
    A = np.zeros_like(precos)  # campo de gauge A_j = (mu_j - r) / sigma_j^2
    for j in range(N_assets):
        A[:, j] = (mu[j] - r) / (sigma[j]**2 + 1e-12) * precos[:, j]

    # Derivadas parciais via diferenças finitas centradas
    dA_dS = np.zeros((N_t - 2, N_assets - 2))
    dA_dt = np.zeros((N_t - 2, N_assets - 2))
    for k in range(1, N_t - 1):
        for j in range(1, N_assets - 1):
            dA_dS[k - 1, j - 1] = (A[k, j + 1] - A[k, j - 1]) / (2 * dS)
            dA_dt[k - 1, j - 1] = (A[k + 1, j] - A[k - 1, j]) / (2 * dt)

    R = dA_dt - dA_dS
    return R

# Verificacao: mu_j = r implica R ~ 0
N_assets, N_t = 5, 500
precos_teste = np.random.lognormal(mean=0, sigma=0.02, size=(N_t, N_assets))
precos_teste = np.cumprod(1 + precos_teste / 100, axis=0) * 100

r_const = 0.05
mu_flat = np.full(N_assets, r_const)   # todos os drifts = r
sigma_test = np.full(N_assets, 0.20)

R_zero = calcular_curvatura(precos_teste, mu_flat, r_const, sigma_test,
                            dt=1/252, dS=1.0)
print(f'Norma de R (mu_j = r): {np.linalg.norm(R_zero):.6e}')
print(f'Max |R|           : {np.max(np.abs(R_zero)):.6e}')
assert np.max(np.abs(R_zero)) < 1e-6, 'R deve ser ~0 quando nao ha arbitragem'
print('OK: R ~ 0 confirmado.')

# Caso com arbitragem: mu != r
mu_arb = np.array([0.05, 0.05, 0.12, 0.05, 0.05])  # ativo 2 com drift anomalo
R_arb = calcular_curvatura(precos_teste, mu_arb, r_const, sigma_test,
                           dt=1/252, dS=1.0)
print(f'\nNorma de R (com arbitragem): {np.linalg.norm(R_arb):.6e}')
print(f'Max |R|                      : {np.max(np.abs(R_arb)):.6e}')
\end{lstlisting}

\begin{ex}
Execute o código acima e verifique que:
\begin{enumerate}
    \item Para \(\mu_j = r\) (todos os ativos), \(\|R\| \approx 0\) dentro da precisão numérica.
    \item Ao introduzir um ativo com \(\mu_j \neq r\), a curvatura torna-se significativamente não nula, sinalizando oportunidade de arbitragem.
\end{enumerate}
\end{ex}

\section{Construção de Estratégias de Arbitragem}
\label{sec:estrategias-arbitragem}

\subsection{Holonomia e Transporte Paralelo}

Quando \(R \neq 0\), o fibrado financeiro possui curvatura não trivial. A holonomia --- o desvio de um vetor após transporte paralelo ao longo de uma curva fechada --- quantifica o ganho de arbitragem. Para uma curva fechada \(\gamma\) no espaço de parâmetros, o operador de holonomia é:
\begin{equation}
    \mathcal{P}\exp\!\left(\oint_\gamma A_\mu\,dx^\mu\right) \neq \mathbb{I},
\end{equation}
onde \(\mathcal{P}\) denota ordenação de caminho (path ordering).

No contexto prático, a estratégia de arbitragem consiste em três etapas:
\begin{enumerate}
    \item \textbf{Encontrar curva fechada}: Identificar um ciclo no grafo de ativos cujo produto de taxas de câmbio (ou preços relativos) desvia de 1.
    \item \textbf{Transporte paralelo}: Propagar um portfólio inicial ao longo da curva usando a conexão de gauge.
    \item \textbf{Extrair lucro}: A diferença entre o portfólio final e o inicial é o lucro de arbitragem.
\end{enumerate}

\begin{lstlisting}[language=Python, caption={Algoritmo de construção de estratégia de arbitragem via holonomia.}]
import numpy as np

def encontrar_ciclo_arbitravel(precos, spreads, limiar=1e-4):
    """
    Identifica ciclo fechado com desvio de não-arbitragem (R != 0).

    Em um mercado de N ativos, um ciclo e uma sequencia de indices
    (i1, i2, ..., ik, i1) onde o produto dos precos relativos se desvia de 1.

    Retorna
    -------
    ciclo : list[int] -- indices do ciclo encontrado
    desvio : float    -- desvio logaritmico (quanto maior, mais lucro)
    """
    N = len(precos)
    # Matriz de precos relativos: X[i,j] = S_i / S_j
    X = precos[:, None] / (precos[None, :] + 1e-12)

    # Para triangulos (i, j, k), verificar:
    # X[i,j] * X[j,k] * X[k,i] deve ser 1. Desvio = produto - 1.
    melhor_desvio = 0.0
    melhor_ciclo = []
    for i in range(N):
        for j in range(i + 1, N):
            for k in range(j + 1, N):
                prod = X[i, j] * X[j, k] * X[k, i]
                desvio = abs(np.log(prod + 1e-12))
                if desvio > melhor_desvio:
                    melhor_desvio = desvio
                    melhor_ciclo = [i, j, k, i]
    return melhor_ciclo, melhor_desvio

def transporte_paralelo(ciclo, A, spreads):
    """
    Realiza transporte paralelo ao longo de um ciclo.

    O vetor de portfolio V e propagado via:
    V_{k+1} = V_k * exp(-A_k * delta_x_k)

    A holonomia e H = V_final / V_inicial - 1 (lucro fracionario).
    """
    V = 1.0
    for idx in range(len(ciclo) - 1):
        i, j = ciclo[idx], ciclo[idx + 1]
        delta_x = spreads[i, j]
        A_ij = A[i, j]
        V *= np.exp(-A_ij * delta_x)
    holonomia = V - 1.0
    return holonomia

def construir_estrategia(precos, mu, r, sigma):
    """
    Constroi estrategia de arbitragem completa.

    Retorna
    -------
    resultado : dict com 'ciclo', 'holonomia', 'lucro_esperado',
                'pesos_portfolio'
    """
    N = len(precos)
    # Spreads de precos
    spreads = np.zeros((N, N))
    for i in range(N):
        for j in range(N):
            spreads[i, j] = precos[j] - precos[i] if i != j else 0.0

    # Campo de gauge
    A = np.zeros((N, N))
    for i in range(N):
        for j in range(N):
            if i != j:
                A[i, j] = (mu[j] - r) / (sigma[j]**2 + 1e-12) - (mu[i] - r) / (sigma[i]**2 + 1e-12)

    ciclo, desvio = encontrar_ciclo_arbitravel(precos, spreads)
    if not ciclo:
        return {'ciclo': [], 'holonomia': 0, 'lucro_esperado': 0}

    hol = transporte_paralelo(ciclo, A, spreads)

    # Pesos: long nos ativos subprecificados, short nos sobreprecificados
    N_cycle = len(ciclo) - 1
    pesos = np.zeros(N)
    for idx in range(N_cycle):
        i, j = ciclo[idx], ciclo[idx + 1]
        sinal = 1.0 if mu[j] > r else -1.0
        pesos[j] += sinal / N_cycle

    lucro_esperado = abs(hol) * 100  # em percentual

    return {
        'ciclo': ciclo,
        'holonomia': hol,
        'lucro_esperado': lucro_esperado,
        'pesos_portfolio': pesos
    }

# Exemplo: 4 ativos com um ativo mal precificado
precos_spot = np.array([100.0, 50.0, 200.0, 75.0])
mu_vec = np.array([0.05, 0.05, 0.15, 0.05])  # ativo 2 tem drift anomalo
r_free = 0.05
sigma_vec = np.array([0.20, 0.25, 0.30, 0.18])

resultado = construir_estrategia(precos_spot, mu_vec, r_free, sigma_vec)
print(f'Ciclo de arbitragem  : {resultado["ciclo"]}')
print(f'Holonomia H          : {resultado["holonomia"]:.6f}')
print(f'Lucro esperado       : {resultado["lucro_esperado"]:.4f}%')
print(f'Pesos do portfolio   : {resultado["pesos_portfolio"]}')
\end{lstlisting}

\begin{defi}[Holonomia financeira]
A holonomia \(H_\gamma = \mathcal{P}\exp(\oint_\gamma A) - 1\) é uma medida direta do lucro de arbitragem disponível ao longo da curva fechada \(\gamma\). Se \(H_\gamma > 0\), existe uma estratégia de investimento com lucro livre de risco e custo inicial nulo.
\end{defi}

\section{Simulação de Mercado com Arbitragem}
\label{sec:mercado-arbitragem}

Nesta seção, construímos um mercado sintético com 3 ativos onde uma inconsistência de precificação gera \(R \neq 0\). Em seguida, implementamos a estratégia de arbitragem e simulamos o P\&L (Profit and Loss) ao longo do tempo.

\begin{lstlisting}[language=Python, caption={Simulação de mercado sintético com 3 ativos e arbitragem. Construção da estratégia e PnL.}]
import numpy as np
import matplotlib.pyplot as plt

# Mercado sintetico: 3 ativos com estrutura de drift inconsistente
# Ativo 0: benchmark (sem arbitragem)
# Ativo 1: benchmark (sem arbitragem)
# Ativo 2: mal precificado (drift acima do justo -> sobrevalorizado)

def simular_mercado_arbitravel(S0_vec, mu_vec, sigma_vec, r, T, N, seed=42):
    """
    Simula mercado com N ativos onde alguns tem drift inconsistente.
    """
    rng = np.random.default_rng(seed)
    dt = T / N
    M = len(S0_vec)
    S = np.zeros((N + 1, M))
    S[0, :] = S0_vec
    for t_idx in range(N):
        eps = rng.normal(0, 1, M)
        # Correlacao parcial via Cholesky
        L = np.array([[1.0, 0.3, 0.1],
                       [0.3, 1.0, 0.2],
                       [0.1, 0.2, 1.0]])
        z = L @ eps
        for j in range(M):
            S[t_idx + 1, j] = S[t_idx, j] * np.exp(
                (mu_vec[j] - 0.5 * sigma_vec[j]**2) * dt
                + sigma_vec[j] * np.sqrt(dt) * z[j]
            )
    return S

# Parametros
S0_vec = np.array([100.0, 100.0, 100.0])
mu_vec = np.array([0.05, 0.05, 0.18])  # ativo 2 com drift anomalo!
sigma_vec = np.array([0.20, 0.20, 0.35])
r_free = 0.05
T_sim, N_sim = 1.0, 252

S_mercado = simular_mercado_arbitravel(S0_vec, mu_vec, sigma_vec,
                                       r_free, T_sim, N_sim)

# Calcular curvatura ao longo do tempo
def curvatura_temporal(S, mu, r, sigma):
    """Calcula norma da curvatura em cada instante de tempo."""
    N_t = S.shape[0]
    R_norm = np.zeros(N_t - 2)
    for t_idx in range(1, N_t - 1):
        A_t = np.array([(mu[j] - r) / (sigma[j]**2 + 1e-12) * S[t_idx, j]
                        for j in range(3)])
        A_prev = np.array([(mu[j] - r) / (sigma[j]**2 + 1e-12) * S[t_idx - 1, j]
                           for j in range(3)])
        A_next = np.array([(mu[j] - r) / (sigma[j]**2 + 1e-12) * S[t_idx + 1, j]
                           for j in range(3)])
        dA_dt = (A_next - A_prev) / (2 * T_sim / N_sim)
        R_norm[t_idx - 1] = np.sqrt(np.sum(dA_dt**2))
    return R_norm

R_norm = curvatura_temporal(S_mercado, mu_vec, r_free, sigma_vec)
print(f'Curvatura media : {np.mean(R_norm):.6e}')
print(f'Curvatura maxima: {np.max(R_norm):.6e}')
print(f'R != 0 confirmado: {np.mean(R_norm) > 1e-6}')

# Construir estrategia e simular PnL
def simular_pnl(S, pesos, dt, r):
    """
    Simula o PnL de uma estrategia de portfolio.
    pesos: ndarray (N_assets,) -- pesos positivos = long, negativos = short
    """
    N_t, N_assets = S.shape
    retornos = np.diff(np.log(S + 1e-12), axis=0)
    pnl = np.zeros(N_t - 1)
    capital = 1.0
    pnl_acum = np.zeros(N_t - 1)
    for t_idx in range(N_t - 1):
        ret_portfolio = np.dot(pesos, retornos[t_idx, :])
        capital *= (1 + ret_portfolio - r * dt)
        pnl_acum[t_idx] = capital - 1.0
    return pnl_acum

# Pesos: long ativo 2 (sobrevalorizado -> short) e long nos outros
pesos_estrategia = np.array([0.5, 0.5, -1.0])
pnl = simular_pnl(S_mercado, pesos_estrategia, T_sim / N_sim, r_free)

# Grafico
fig, (ax1, ax2) = plt.subplots(2, 1, figsize=(10, 8), sharex=True)

ax1.plot(S_mercado[:, 0], label='Ativo 0 (benchmark)', lw=1.2)
ax1.plot(S_mercado[:, 1], label='Ativo 1 (benchmark)', lw=1.2)
ax1.plot(S_mercado[:, 2], label='Ativo 2 (mal precificado)', lw=1.5)
ax1.set_ylabel('Preco')
ax1.set_title('Mercado Sintetico com Arbitragem')
ax1.legend()
ax1.grid(True, alpha=0.3)

t_dias = np.arange(len(pnl))
ax2.plot(t_dias, pnl * 100, 'g-', lw=1.5)
ax2.fill_between(t_dias, 0, pnl * 100, alpha=0.15, color='green')
ax2.axhline(0, color='gray', ls='--')
ax2.set_xlabel('Dias uteis')
ax2.set_ylabel('PnL Acumulado (%)')
ax2.set_title('PnL da Estrategia de Arbitragem')
ax2.grid(True, alpha=0.3)

plt.tight_layout()
plt.savefig('mercado_arbitragem_pnl.pdf')
plt.show()

print(f'\nPnL final: {pnl[-1] * 100:.2f}%')
print(f'Sharpe ratio: {np.mean(pnl) / (np.std(pnl) + 1e-12) * np.sqrt(252):.2f}')
\end{lstlisting}

\begin{ex}
Analise os resultados da simulação:
\begin{enumerate}
    \item O ativo 2, com \(\mu = 0.18\) quando deveria ter \(\mu = r = 0.05\), apresenta trajetória consistentemente superior --- evidência de má precificação.
    \item A curvatura \(R\) é não nula em todo o período, confirmando a presença de oportunidade de arbitragem.
    \item A estratégia \textit{long-short} (comprado nos ativos corretamente precificados, vendido no ativo sobrevalorizado) gera PnL positivo e crescente.
    \item O Sharpe ratio elevado indica que o lucro não é compensação por risco, mas sim exploração de ineficiência de mercado --- a assinatura da arbitragem.
\end{enumerate}
\end{ex}

\section{Visualização da Curvatura}
\label{sec:visualizacao-curvatura}

A visualização da curvatura como função do tempo e da composição da carteira permite identificar regiões de alta curvatura --- os ``hotspots'' de arbitragem. Nesta seção, geramos gráficos de calor (\textit{heatmaps}) e superfícies 3D.

\begin{lstlisting}[language=Python, caption={Visualização da curvatura: heatmap temporal e superfície 3D em função da carteira.}]
import numpy as np
import matplotlib.pyplot as plt
from matplotlib import cm

# Gerar grade de portfolios (composicoes)
def grade_portfolios(N_assets, n_pts=50):
    """Gera grade de pesos para visualizacao (simplex para 3 ativos)."""
    if N_assets == 3:
        x = np.linspace(0, 1, n_pts)
        y = np.linspace(0, 1, n_pts)
        X, Y = np.meshgrid(x, y)
        # Restricao: w0 + w1 + w2 = 1, todos >= 0
        mask = X + Y <= 1.0
        return X, Y, mask
    else:
        raise ValueError('Visualizacao implementada para 3 ativos')

def curvatura_portfolio(mu, r, sigma, w):
    """
    Calcula curvatura para uma carteira com pesos w.

    R_portfolio = || (mu - r*1) x w || / (w^T Sigma w)
    onde x denota produto vetorial (assimetria de Sharpe).
    """
    excesso = mu - r
    # Componente antissimetrica: excesso_i * w_j - excesso_j * w_i
    R_val = 0.0
    for i in range(len(w)):
        for j in range(len(w)):
            R_val += (excesso[i] * w[j] - excesso[j] * w[i])**2
    R_val = np.sqrt(R_val)
    return R_val

# Parametros
mu_vec = np.array([0.05, 0.07, 0.14])
sigma_vec = np.array([0.20, 0.25, 0.30])
r_free = 0.05

X, Y, mask = grade_portfolios(3, n_pts=80)
Z = np.zeros_like(X)
for i in range(X.shape[0]):
    for j in range(X.shape[1]):
        if mask[i, j]:
            w0 = X[i, j]
            w1 = Y[i, j]
            w2 = 1.0 - w0 - w1
            w = np.array([w0, w1, w2])
            Z[i, j] = curvatura_portfolio(mu_vec, r_free, sigma_vec, w)
        else:
            Z[i, j] = np.nan

# Heatmap da curvatura no simplex
fig, (ax1, ax2) = plt.subplots(1, 2, figsize=(14, 5.5))

im = ax1.pcolormesh(X, Y, Z, shading='auto', cmap='inferno')
ax1.set_xlabel('Peso Ativo 0')
ax1.set_ylabel('Peso Ativo 1')
ax1.set_title('Curvatura $R$ no Simplex de Carteiras')
cbar = fig.colorbar(im, ax=ax1, label='$||R||$')
# Linha de triangulo
ax1.plot([0, 1, 0, 0], [0, 0, 1, 0], 'w-', lw=1.5, alpha=0.5)

# Curvatura temporal (simulacao)
N_dias = 252
t_grid = np.linspace(0, 1, N_dias)
R_temporal = np.zeros(N_dias)
for k in range(N_dias):
    # Simular precos com drift variavel
    fator = 1 + 0.3 * np.sin(2 * np.pi * k / N_dias)  # sazonalidade
    mu_t = mu_vec * fator
    excesso_t = np.linalg.norm(mu_t - r_free)
    R_temporal[k] = excesso_t

ax2.plot(t_grid, R_temporal, 'b-', lw=1.5)
ax2.fill_between(t_grid, 0, R_temporal, alpha=0.2, color='blue')
ax2.axhline(np.mean(R_temporal), color='red', ls='--',
            label=f'Media = {np.mean(R_temporal):.4f}')
ax2.set_xlabel('Tempo (fracao do ano)')
ax2.set_ylabel('$||R||$')
ax2.set_title('Curvatura ao Longo do Tempo')
ax2.legend()
ax2.grid(True, alpha=0.3)

plt.tight_layout()
plt.savefig('curvatura_visualizacao.pdf')
plt.show()

# Identificar regioes de alta curvatura
limiar_alto = np.nanpercentile(Z, 90)
indices_altos = np.argwhere((Z >= limiar_alto) & ~np.isnan(Z))
print(f'Regioes de alta curvatura (>{limiar_alto:.4f}): {len(indices_altos)} pontos')
print(f'Curvatura maxima: {np.nanmax(Z):.4f}')
print(f'Curvatura minima: {np.nanmin(Z):.4f}')
\end{lstlisting}

\begin{ex}
Interpretação dos gráficos:
\begin{enumerate}
    \item O \textbf{heatmap no simplex} mostra que carteiras concentradas no ativo com maior desvio de precificação (Ativo 2, canto inferior direito) apresentam curvatura máxima --- estas são as carteiras com maior potencial de arbitragem.
    \item A \textbf{série temporal} revela que a curvatura não é constante: oscila com a dinâmica do mercado, criando janelas de oportunidade que o arbitrador deve identificar e explorar.
    \item Regiões de curvatura acima do percentil 90 são candidatas prioritárias para execução de estratégias.
\end{enumerate}
\end{ex}

\section{Exemplo Completo: Mini-Framework GAT}
\label{sec:framework-gat}

Esta seção consolida todos os conceitos em um mini-framework orientado a objetos, com classes para os componentes fundamentais da GAT: \texttt{Gauge}, \texttt{MarketBundle}, \texttt{Connection} e \texttt{Curvature}. O framework é autocontido e funcional, permitindo ao estudante instanciar mercados, calcular curvaturas e construir estratégias de arbitragem.

\begin{lstlisting}[language=Python, caption={Mini-framework GAT: classes para Gauge, MarketBundle, Connection e Curvature. Código funcional completo.}]
import numpy as np
from dataclasses import dataclass, field
from typing import List, Tuple, Optional

# =====================================================================
# 1. GAUGE: encapsula o campo de gauge A_mu
# =====================================================================
@dataclass
class Gauge:
    """Campo de gauge A_mu no fibrado financeiro.

    Atributos
    ---------
    valores : ndarray (N_assets, N_assets) -- matriz do campo de gauge
    nome    : str -- identificador do gauge
    """
    valores: np.ndarray
    nome: str = 'A'

    @classmethod
    def from_drifts(cls, mu, r, sigma, precos):
        """Constroi gauge a partir de drifts e precos spot."""
        N = len(mu)
        A = np.zeros((N, N))
        for i in range(N):
            for j in range(N):
                if i != j:
                    excesso_i = (mu[i] - r) / (sigma[i]**2 + 1e-12)
                    excesso_j = (mu[j] - r) / (sigma[j]**2 + 1e-12)
                    A[i, j] = (excesso_j * precos[j] - excesso_i * precos[i])
        return cls(valores=A, nome='A')

    def __repr__(self):
        return f'Gauge({self.nome}, shape={self.valores.shape})'


# =====================================================================
# 2. MARKET BUNDLE: representa o fibrado de mercado
# =====================================================================
@dataclass
class MarketBundle:
    """Fibrado de mercado: espaco base (tempo) x fibra (precos).

    Atributos
    ---------
    precos     : ndarray (N_t, N_assets)
    tempos     : ndarray (N_t,)
    nomes      : list[str]
    taxa_livre : float
    """
    precos: np.ndarray
    tempos: np.ndarray
    nomes: List[str]
    taxa_livre: float = 0.05

    @property
    def n_assets(self):
        return self.precos.shape[1]

    @property
    def n_periodos(self):
        return self.precos.shape[0]

    def retornos_log(self):
        return np.diff(np.log(self.precos + 1e-12), axis=0)

    def __repr__(self):
        return (f'MarketBundle({self.n_assets} ativos, '
                f'{self.n_periodos} periodos, r={self.taxa_livre})')


# =====================================================================
# 3. CONNECTION: conexao de gauge no fibrado
# =====================================================================
@dataclass
class Connection:
    """Conexao de gauge no fibrado financeiro.

    Define transporte paralelo e derivada covariante.
    """
    gauge: Gauge
    bundle: MarketBundle
    dt: float = 1 / 252

    def transporte_paralelo(self, vetor, passo):
        """
        Transporta vetor ao longo da fibra usando a conexao.

        V(t+dt) = V(t) - A_mu * V(t) * dx^mu
        """
        A = self.gauge.valores
        delta_v = -A @ vetor * self.dt
        return vetor + delta_v

    def holonomia_ciclo(self, ciclo):
        """
        Calcula holonomia ao longo de um ciclo fechado.

        H = P exp(oint A) - I
        """
        V = np.eye(self.bundle.n_assets)
        A = self.gauge.valores
        for idx in range(len(ciclo) - 1):
            i, j = ciclo[idx], ciclo[idx + 1]
            V = V @ (np.eye(self.bundle.n_assets) - A * (self.bundle.precos[-1, j] - self.bundle.precos[-1, i]) / self.bundle.precos[-1, i])
        holonomia = np.trace(V) / self.bundle.n_assets - 1.0
        return holonomia

    def derivada_covariante(self, secao):
        """
        Derivada covariante de uma secao do fibrado.

        D_mu psi = partial_mu psi + A_mu psi
        """
        A = self.gauge.valores
        return secao + A @ secao * self.dt

    def __repr__(self):
        return f'Connection(gauge={self.gauge.nome}, bundle={self.bundle})'


# =====================================================================
# 4. CURVATURE: tensor de curvatura e analise de arbitragem
# =====================================================================
@dataclass
class Curvature:
    """Tensor de curvatura R e ferramentas de analise.

    R_{mu,nu} = partial_mu A_nu - partial_nu A_mu + [A_mu, A_nu]
    """
    connection: Connection
    _R: Optional[np.ndarray] = field(default=None, repr=False)
    _calculada: bool = field(default=False, repr=False)

    def calcular(self):
        """Calcula o tensor de curvatura numericamente."""
        A = self.connection.gauge.valores
        N = A.shape[0]
        dt = self.connection.dt

        # Aproximacao: R_{ij} = (A_{ij}(t+dt) - A_{ij}(t-dt)) / (2*dt) + [A, A]_{ij}
        R = np.zeros((N, N))
        comutador = A @ A - A @ A  # zero para matrizes (abeliano)
        # Para capturar nao-abeliano, usariamos multiplos gauges

        # Componente de derivada temporal
        precos = self.connection.bundle.precos
        if precos.shape[0] >= 3:
            for i in range(N):
                for j in range(N):
                    dS_i = (precos[-1, i] - precos[0, i]) / (precos.shape[0] * dt + 1e-12)
                    dS_j = (precos[-1, j] - precos[0, j]) / (precos.shape[0] * dt + 1e-12)
                    R[i, j] = (A[i, j] * dS_j - A[j, i] * dS_i) / (precos[-1, i] + 1e-12)

        self._R = R
        self._calculada = True
        return R

    @property
    def norma(self):
        """Norma de Frobenius da curvatura."""
        if not self._calculada:
            self.calcular()
        return np.linalg.norm(self._R, 'fro')

    @property
    def tem_arbitragem(self):
        """Verifica se ha arbitragem (R significativamente nao nulo)."""
        return self.norma > 1e-6

    def decompor(self):
        """
        Decompoe curvatura em componentes:
        - pura (abeliana): dA
        - mista (nao-abeliana): [A, A]
        """
        if not self._calculada:
            self.calcular()
        A = self.connection.gauge.valores
        dA = self._R  # componente abeliana
        comutador = A @ A - A @ A  # zero para gauge unico
        return {'dA': dA, 'comutador': comutador}

    def __repr__(self):
        status = 'calculada' if self._calculada else 'nao calculada'
        return (f'Curvature(status={status}, ||R||={self.norma:.6e}, '
                f'arbitragem={self.tem_arbitragem})')


# =====================================================================
# 5. ANALISADOR DE ARBITRAGEM: orquestrador do framework
# =====================================================================
class AnalisadorArbitragem:
    """Orquestrador do mini-framework GAT.

    Integra MarketBundle -> Gauge -> Connection -> Curvature
    e gera relatorio completo de arbitragem.
    """

    def __init__(self, precos, mu, sigma, r, nomes=None, dt=1/252):
        self.mu = np.asarray(mu)
        self.sigma = np.asarray(sigma)
        self.r = r
        self.dt = dt

        N_t, N_assets = precos.shape
        t_grid = np.linspace(0, N_t * dt, N_t)
        if nomes is None:
            nomes = [f'Ativo_{i}' for i in range(N_assets)]

        self.bundle = MarketBundle(precos=precos, tempos=t_grid,
                                   nomes=nomes, taxa_livre=r)
        self.gauge = Gauge.from_drifts(mu, r, sigma, precos[-1, :])
        self.connection = Connection(gauge=self.gauge, bundle=self.bundle, dt=dt)
        self.curvature = Curvature(connection=self.connection)

    def analisar(self):
        """Executa analise completa de arbitragem."""
        R = self.curvature.calcular()
        norma_R = self.curvature.norma

        # Encontrar melhor ciclo de arbitragem
        N = self.bundle.n_assets
        melhor_ciclo, melhor_lucro = [], 0.0

        precos_spot = self.bundle.precos[-1, :]
        for i in range(N):
            for j in range(i + 1, N):
                for k in range(j + 1, N):
                    # Triangulo i -> j -> k -> i
                    razao = (precos_spot[j] / precos_spot[i]
                             * precos_spot[k] / precos_spot[j]
                             * precos_spot[i] / precos_spot[k])
                    lucro = abs(1.0 - razao)
                    if lucro > melhor_lucro:
                        melhor_lucro = lucro
                        melhor_ciclo = [i, j, k, i]

        return {
            'norma_R': norma_R,
            'tem_arbitragem': self.curvature.tem_arbitragem,
            'n_ativos': N,
            'melhor_ciclo': melhor_ciclo,
            'lucro_triangular': melhor_lucro,
            'gauge': self.gauge,
            'curvature': self.curvature,
        }

    def relatorio(self):
        """Gera relatorio formatado."""
        analise = self.analisar()
        print('=' * 60)
        print('  RELATORIO DE ANALISE GAT')
        print('=' * 60)
        print(f'  Ativos analisados      : {analise["n_ativos"]}')
        print(f'  Taxa livre de risco    : {self.r:.4f}')
        print(f'  Norma ||R||            : {analise["norma_R"]:.6e}')
        print(f'  Arbitragem detectada   : {"SIM" if analise["tem_arbitragem"] else "NAO"}')
        print(f'  Melhor ciclo           : {analise["melhor_ciclo"]}')
        print(f'  Lucro triangular       : {analise["lucro_triangular"]*100:.4f}%')
        print('=' * 60)

        # Diagnostico por ativo
        print('\n  Diagnostico por ativo:')
        print(f'  {"Ativo":<12} {"Drift":>8} {"Sigma":>8} {"Excesso":>10} {"Sharpe":>8}')
        print('  ' + '-' * 50)
        for j in range(analise['n_ativos']):
            excesso = self.mu[j] - self.r
            sharpe = excesso / (self.sigma[j] + 1e-12)
            nome = self.bundle.nomes[j]
            print(f'  {nome:<12} {self.mu[j]:>8.4f} {self.sigma[j]:>8.4f} '
                  f'{excesso:>10.4f} {sharpe:>8.2f}')
        print()

        return analise


# =====================================================================
# 6. EXECUCAO DO EXEMPLO COMPLETO
# =====================================================================
if __name__ == '__main__':
    # Dados de entrada
    N_assets = 4
    N_periodos = 500
    r_free = 0.05
    dt = 1 / 252

    # Drifts: ativo 2 tem drift anomalo (arbitragem)
    mu = np.array([0.05, 0.05, 0.16, 0.05])
    sigma = np.array([0.20, 0.22, 0.35, 0.18])
    nomes = ['PETR4', 'VALE3', 'ARB1', 'ITUB4']
    S0 = np.array([35.0, 68.0, 100.0, 32.0])

    # Simular precos
    rng = np.random.default_rng(42)
    precos = np.zeros((N_periodos, N_assets))
    precos[0, :] = S0
    for t in range(N_periodos - 1):
        eps = rng.normal(0, 1, N_assets)
        for j in range(N_assets):
            precos[t + 1, j] = precos[t, j] * np.exp(
                (mu[j] - 0.5 * sigma[j]**2) * dt
                + sigma[j] * np.sqrt(dt) * eps[j]
            )

    # Instanciar e executar o framework
    analisador = AnalisadorArbitragem(precos, mu, sigma, r_free, nomes, dt)
    resultado = analisador.relatorio()

    # Verificacao adicional
    print('\nDecomposicao da curvatura:')
    decomp = analisador.curvature.decompor()
    print(f'  ||dA||        = {np.linalg.norm(decomp["dA"], "fro"):.6e}')
    print(f'  ||[A,A]||     = {np.linalg.norm(decomp["comutador"], "fro"):.6e}')
    print(f'  Interpretacao : componente abeliana domina (gauge unico)')
\end{lstlisting}

\begin{defi}[Framework GAT mínimo]
O mini-framework acima implementa o ciclo completo da GAT:
\begin{enumerate}
    \item \textbf{Gauge.from\_drifts}: Constrói o campo de gauge a partir de dados observáveis de mercado (preços, drifts, volatilidades).
    \item \textbf{MarketBundle}: Encapsula a geometria do fibrado --- base temporal e fibra de preços.
    \item \textbf{Connection}: Define transporte paralelo e derivada covariante, operações fundamentais para propagação de carteiras.
    \item \textbf{Curvature}: Calcula o tensor \(R\) e verifica a condição de integrabilidade (ausência de arbitragem \(\iff R = 0\)).
    \item \textbf{AnalisadorArbitragem}: Orquestrador que integra todos os componentes e gera diagnóstico completo.
\end{enumerate}
\end{defi}

\begin{ex}
Modifique o exemplo acima para:
\begin{enumerate}
    \item Adicionar um quinto ativo com drift ainda maior (\(\mu = 0.25\)) e verificar o aumento na norma de \(R\).
    \item Introduzir correlação não trivial entre os ativos (matriz de covariância não diagonal) e observar o efeito na curvatura.
    \item Implementar o caso não-abeliano com dois gauges \(A^{(1)}\) e \(A^{(2)}\) e calcular o comutador \([A^{(1)}, A^{(2)}]\).
    \item Estender o \texttt{AnalisadorArbitragem} para buscar ciclos de comprimento arbitrário (não apenas triângulos) usando o algoritmo de Bellman-Ford para detecção de ciclos negativos.
\end{enumerate}
\end{ex}

\section{Considerações Finais do Capítulo}
\label{sec:consideracoes-finais-comp}

Este capítulo demonstrou a viabilidade computacional da Teoria Geométrica de Arbitragem. Os principais resultados são:

\begin{enumerate}
    \item A simulação de GBM fornece a base estocástica para testar a GAT em ambientes controlados.
    \item O cálculo numérico da curvatura confirma a equivalência fundamental: \(R = 0 \iff\) ausência de arbitragem.
    \item A holonomia fornece um algoritmo construtivo para extrair lucros de arbitragem a partir de \(R \neq 0\).
    \item A simulação de mercado sintético valida a estratégia com PnL positivo e Sharpe elevado.
    \item A visualização da curvatura revela \textit{hotspots} de oportunidade no espaço de carteiras e no tempo.
    \item O mini-framework GAT consolida todos os conceitos em código reutilizável e extensível.
\end{enumerate}

O estudante deve agora ser capaz de implementar, testar e estender a GAT para cenários mais complexos, incluindo múltiplos fatores de risco, custos de transação e mercados incompletos --- tópicos abordados nos capítulos avançados seguintes.


